\documentclass[12pt,a4paper]{article}
\usepackage[utf8]{inputenc}
\usepackage[ukrainian]{babel}
\usepackage{geometry}
\geometry{left=25mm,right=15mm,top=20mm,bottom=20mm}
\usepackage{hyperref}
\usepackage{tabularx}
\usepackage{booktabs}
\usepackage{enumitem}
\usepackage{titlesec}
\usepackage{xcolor}

\hypersetup{
    colorlinks=true,
    linkcolor=blue,
    filecolor=magenta,      
    urlcolor=cyan,
    pdftitle={Специфікація вимог до ПЗ (SRS) - web-app_for_design-references},
    pdfpagemode=FullScreen,
}

\title{\textbf{Специфікація вимог до програмного забезпечення (SRS)}\\
\large Система каталогізації та швидкого пошуку дизайн-референсів\\
\normalsize Версія 1.0 (Базова лінія SPEC-GATE Approved)}

\author{Андрій Волошин (гр. КН-31)\\
КПІ ім. Ігоря Сікорського\\
Дисципліна: Проектування інформаційних систем\\
Кодинговий агент: Antigravity AI Coding Assistant (Gemini 3.8 Flash)}
\date{3 жовтня 2026 р.}

\begin{document}

\maketitle
\tableofcontents
\newpage

\section{Вступ (Introduction)}

\subsection{Призначення документа}
Цей документ містить повний, несуперечливий та формалізований опис вимог до програмного забезпечення \textbf{«Система каталогізації та швидкого пошуку дизайн-референсів»} (\texttt{web-app\_for\_design-references}). Документ розроблено відповідно до рекомендацій стандарту IEEE 830 та навчальної програми дисципліни «Проектування інформаційних систем» у межах методології Specification-Driven Development (SDD). Документ виступає обов'язковим еталоном для декомпозиції задач, розробки тестів за методикою TDD/BDD, імплементації вихідного коду та проведення перевірок на контрольних точках.

\subsection{Межі проєкту (Scope)}
Програмний продукт є персоналізованим інструментом для дизайнерів цифрових продуктів (UI/UX, веб, мобільні інтерфейси), який вирішує проблему швидкого збереження візуальних матеріалів, їх категоризації, миттєвої фільтрації та пошуку без перевантаження сторінок і без надмірних форм введення даних (нульове тертя, Zero-Friction UX).

\subsection{Визначення, акроніми та скорочення}
\begin{itemize}[noitemsep]
    \item \textbf{SRS} --- Software Requirements Specification (Специфікація вимог до програмного забезпечення).
    \item \textbf{SDD} --- Specification-Driven Development (Розробка, керована специфікацією).
    \item \textbf{SKED} --- Structured Knowledge Elicitation Dialogue (Структурований діалог виявлення знань).
    \item \textbf{SPEC-GATE} --- Контрольна точка верифікації та офіційного затвердження специфікації.
    \item \textbf{BDD} --- Behavior-Driven Development (Розробка, керована поведінкою).
    \item \textbf{AC} --- Acceptance Criteria (Критерії прийняття вимог).
    \item \textbf{Референс (Reference)} --- Одиниця даних із зображенням, назвою, джерелом (URL), категорією, набором тегів та нотатками.
    \item \textbf{Фасетний тег (Facet Tag)} --- Тег формату \texttt{namespace:value} для параметризованої фільтрації.
\end{itemize}

\subsection{Посилання на пов'язані документи}
\begin{enumerate}[noitemsep]
    \item IEEE Std 830-1998: IEEE Recommended Practice for Software Requirements Specifications.
    \item Регламент контролю версій: \texttt{docs/git-workflow.md}.
    \item Протокол SKED-діалогу: \texttt{logs/sked-summary.md}.
    \item Набір поведінкових BDD-сценаріїв: \texttt{spec/bdd-scenarios.md}.
    \item Початкова матриця простежуваності: \texttt{spec/traceability-matrix.md}.
\end{enumerate}

\section{Загальний опис системи (Overall Description)}

\subsection{Контекст продукту}
Система функціонує як легковажний автономний веб-застосунок (Single Page Application) з локальним зберіганням стану. Застосунок не залежить від сторонніх платних хмарних сервісів, надаючи дизайнеру миттєвий робочий простір для організації натхнення безпосередньо у процесі проєктування інтерфейсів.

\subsection{Класи користувачів}
\begin{itemize}
    \item \textbf{UI/UX Дизайнер (Основний користувач):} Потребує швидкого збереження референсів (до 2--3 дій), моментальної фільтрації за стилями та категоріями під час пошуку ідей.
    \item \textbf{Арт-директор / Лід команди:} Використовує систему для перегляду тематичних колекцій та презентації концепцій замовникам.
\end{itemize}

\subsection{Операційне середовище}
Клієнтська частина функціонує у сучасних веб-браузерах (Google Chrome, Firefox, Safari, Edge). Серверна/модельна частина базується на середовищі Python 3.10+ або нативному рушії веб-середовища.

\subsection{Загальні обмеження та припущення}
\begin{itemize}[noitemsep]
    \item Принцип нульового тертя (Zero-Friction UX): мінімум обов'язкових полів.
    \item Висока швидкодія: виконання фільтрації локально без мережевих затримок.
    \item Конвенція версіювання: зміни вимог фіксуються комітами з префіксом \texttt{spec:}.
    \item Каталог містить до 5000 елементів у межах релізу v1.0, що гарантує миттєву обробку в оперативній пам'яті.
\end{itemize}

\section{Межі системи (System Scope: In-Scope vs Out-of-Scope)}

\subsection{Функціональність релізу v1.0 (In-Scope)}
\begin{enumerate}[noitemsep]
    \item Створення референсу за посиланням (URL), назвою та категорією.
    \item Прив'язка до єдиної взаємовиключної категорії (Primary Category) з дефолтом \texttt{General}.
    \item Організація референсів у колекції (N:M).
    \item Додавання та видалення нормалізованих тегів і фасетів формату \texttt{namespace:value}.
    \item Миттєва кон'юнктивна фільтрація за декількома тегами (логічне AND).
    \item Параметризована фільтрація за фасетами (наприклад, \texttt{style:minimalism}).
    \item Повнотекстовий нечутливий до регістру пошук за назвою, нотатками та тегами.
    \item Захист цілісності: збереження референсів при видаленні категорії (автоматичне перепризначення в \texttt{General}).
    \item Експорт та імпорт даних у форматі JSON.
\end{enumerate}

\subsection{Функціональність поза межами v1.0 (Out-of-Scope)}
\begin{itemize}[noitemsep]
    \item \textbf{Computer Vision (ШІ-зір):} Автоматичне нейромережеве розпізнавання палітри та стилю (винесено в майбутній адендум).
    \item \textbf{Багатокористувацька онлайн-співпраця:} Real-time синхронізація через WebSockets.
    \item \textbf{Соціальний маркетплейс:} Публічні рейтинги, лайки, продаж дизайн-робіт.
    \item \textbf{Платіжні шлюзи та платні плани:} Не передбачені у навчальному проєкті.
\end{itemize}

\section{Функціональні вимоги (Functional Requirements)}

\begin{description}
    \item[\textbf{REQ-F-001: Створення референсу із зовнішнього джерела}] \hfill \\
    Система повинна забезпечувати створення нового референсу за вказаними назвою, URL-адресою джерела, необов'язковим прев'ю, категорією, набором тегів та нотатками.
    \begin{itemize}[noitemsep]
        \item \textbf{AC-F-001.1:} При непорожній назві та валідному URL референс успішно зберігається.
        \item \textbf{AC-F-001.2:} При порожній назві або некоректному URL виводиться повідомлення про помилку валідації.
        \item \textbf{AC-F-001.3:} Заборонено збереження адрес із небезпечними протоколами (не \texttt{http://} і не \texttt{https://}).
    \end{itemize}

    \item[\textbf{REQ-F-002: Асоціація референсу з категорією}] \hfill \\
    Кожен референс строго асоціюється з однією категорією (відношення 1:1).
    \begin{itemize}[noitemsep]
        \item \textbf{AC-F-002.1:} Користувач може призначити одну з існуючих категорій.
        \item \textbf{AC-F-002.2:} Якщо категорія не вказана, системно встановлюється \texttt{General}.
    \end{itemize}

    \item[\textbf{REQ-F-003: Групування референсів у колекції (N:M)}] \hfill \\
    Система повинна підтримувати додавання референсу до кількох користувацьких колекцій.
    \begin{itemize}[noitemsep]
        \item \textbf{AC-F-003.1:} Референс може входити в довільну кількість колекцій.
        \item \textbf{AC-F-003.2:} Видалення з колекції не видаляє референс із системи.
    \end{itemize}

    \item[\textbf{REQ-F-004: Управління тегами та фасетами}] \hfill \\
    Система дозволяє додавати нормалізовані теги та фасети \texttt{namespace:value}.
    \begin{itemize}[noitemsep]
        \item \textbf{AC-F-004.1:} Теги нормалізуються до нижнього регістру без зайвих пробілів.
        \item \textbf{AC-F-004.2:} Теги з двокрапкою розпізнаються як префіксні фасети.
    \end{itemize}

    \item[\textbf{REQ-F-005: Миттєва кон'юнктивна фільтрація}] \hfill \\
    Фільтрація списку матеріалів за довільною комбінацією тегів.
    \begin{itemize}[noitemsep]
        \item \textbf{AC-F-005.1:} Референс відображається у вибірці тільки за умови наявності всіх обраних тегів (AND-семантика).
    \end{itemize}

    \item[\textbf{REQ-F-006: Повнотекстовий пошук}] \hfill \\
    Пошук за входженням підрядка в назву, нотатки та теги без урахування регістру.
    \begin{itemize}[noitemsep]
        \item \textbf{AC-F-006.1:} Збіг у будь-якому з текстових полів повертає референс у результатах.
    \end{itemize}

    \item[\textbf{REQ-F-007: Перегляд детальної картки}] \hfill \\
    Відображення повної інформації про вибраний об'єкт.
    \begin{itemize}[noitemsep]
        \item \textbf{AC-F-007.1:} Картка відображає прев'ю, клікабельне посилання, категорію, теги та нотатки.
    \end{itemize}

    \item[\textbf{REQ-F-008: Редагування та видалення}] \hfill \\
    Оновлення даних або повне видалення референсу з каталогу.
    \begin{itemize}[noitemsep]
        \item \textbf{AC-F-008.1:} Зміни миттєво зберігаються у сховищі.
        \item \textbf{AC-F-008.2:} Видалення очищує референс з усіх колекцій без порушення цілісності.
    \end{itemize}
\end{description}

\section{Нефункціональні вимоги (Non-Functional Requirements)}

\begin{description}
    \item[\textbf{REQ-NF-001: Швидкодія (Performance)}] \hfill \\
    \textbf{AC-NF-001.1:} Для каталогу обсягом до 1000 елементів фільтрація за 3 тегами або повнотекстовий пошук виконуються за час \textbf{$<$ 50 мілісекунд}.

    \item[\textbf{REQ-NF-002: Зручність використання (Usability)}] \hfill \\
    \textbf{AC-NF-002.1:} Додавання нового референсу із зазначенням URL, категорії та 3 тегів здійснюється не більше ніж за \textbf{3 дії (кліки)}.

    \item[\textbf{REQ-NF-003: Надійність та цілісність (Reliability)}] \hfill \\
    \textbf{AC-NF-003.1:} Видалення категорії не видаляє асоційовані референси, а переводить їх у категорію \texttt{General}. Втрата референсів заборонена (0\% data loss).

    \item[\textbf{REQ-NF-004: Безпека (Security)}] \hfill \\
    \textbf{AC-NF-004.1:} Дозволені лише схеми \texttt{http://} та \texttt{https://}. Схеми \texttt{javascript:} та \texttt{data:} блокуються. Текстові поля санітизуються проти XSS.

    \item[\textbf{REQ-NF-005: Супровідність (Maintainability)}] \hfill \\
    \textbf{AC-NF-005.1:} Чітке відокремлення бізнес-логіки. Покриття модульних тестів (Unit/BDD) становить не менше 95\%.
\end{description}

\section{Вимоги до інтерфейсів (Interface Requirements)}
\subsection{Інтерфейс користувача}
Адаптивна сітка карток референсів (Masonry / Grid layout) із зображеннями-прев'ю, панель швидкого пошуку із живим оновленням результатів та бічна/верхня панель категорій із чіпсами тегів.

\subsection{Програмний інтерфейс (Data Interface)}
Модуль моделей надає уніфіковані методи: \texttt{add\_reference}, \texttt{filter\_by\_tags}, \texttt{search}, \texttt{remove\_category}. Серіалізація стану здійснюється за стандартизованою схемою JSON.

\section{Контрольна точка SPEC-GATE}
Ця специфікація пройшла офіційну перевірку інженера на контрольній точці \textbf{SPEC-GATE} (протокол у \texttt{logs/spec-gate.md}) та затверджена зі статусом \textbf{APPROVED}. Вона слугує незмінною базовою лінією (Baseline v1.0). Подальші модифікації допускаються лише через процедуру Адендуму (Addendum).

\end{document}
